%========================================
% 8 | Saṅghānussati
%========================================
% title-en: Recollection of the Sangha
\section*{e8 | Saṅghānussati}

\begin{chantsubsection}
  \chantnote
    {\{Handa mayaṃ saṅghānussatinayaṃ karomase\}}
  \chantnotetrans
    {\{Now let us chant in recollection of the Saṅgha.\}}
\end{chantsubsection}

\begin{chantsubsection}
  \chantline
    {Supaṭipanno bhagavato sāvakasaṅgho}
    {The Saṅgha of the Blessed One’s disciples has practiced well.}

  \chantline
    {Ujupaṭipanno bhagavato sāvakasaṅgho}
    {The Saṅgha of the Blessed One’s disciples has practiced uprightly.}

  \chantline
    {Ñāyapaṭipanno bhagavato sāvakasaṅgho}
    {The Saṅgha of the Blessed One’s disciples has practiced rightly.}

  \chantline
    {Sāmīcipaṭipanno bhagavato sāvakasaṅgho}
    {The Saṅgha of the Blessed One’s disciples has practiced properly.}

  \chantline
    {Yadidaṃ cattāri purisayugāni aṭṭha purisapuggalā}
    {That is to say, the four pairs of persons, the eight kinds of individuals.}

  \chantline
    {Esa bhagavato sāvakasaṅgho}
    {This is the Saṅgha of the Blessed One’s disciples.}

  \chantline
    {Āhuneyyo pāhuneyyo dakkhiṇeyyo añjalīkaraṇīyo}
    {Worthy of gifts, worthy of hospitality, worthy of offerings, worthy of reverential salutation.}

  \chantline
    {Anuttaraṃ puññakkhettaṃ lokassāti}
    {The unsurpassed field of merit for the world.}
\end{chantsubsection}
